\section{Discussion}
\subsection{Explanation of key findings}
The overall result reveals a two-sided picture when the data structure is examined more closely. The advantage of GPT-4o-mini is real at the aggregate level (+4.6 points, exceeding 4,999/5,000 random draws), but it is concentrated and uneven: 63\% of the additional killed mutants come from the 3 largest parser groups in JacksonCore. Classes such as UTF8StreamJsonParser have dense state logic and many conditional branches - an environment where the LLM's scoring (1-5) has enough room to distinguish between shallow mutants (small constant changes with little impact) and deep mutants (those that break boundary conditions in parsing loops). Conversely, in Codec, the pool's baseline mutation score was 84.9\% - a clear ceiling effect: when the test suite already killed almost all mutants, "skillful selection" and "random selection" converge to the same result, accurately reflected in the ratio of 4 wins/4 losses/5 ties and $r\_rb=+0.028 \approx 0$.

The gap between the two statistical conclusions (group-weighted p=0.059 vs. mutant-weighted p=0.0002) is a mechanical consequence of the design: group-level testing assigns equal weight to a group with k=1 and a group with k=173, and 30\% of pairs are tied (mainly small groups, coarse MS quantization, plus the ceiling effect in Codec), which negates the Wilcoxon test's ranking information. This is not a contradiction; it reflects two different questions: "Does GPT win in most classes?" (not yet proven) and "Does the mutant set selected by GPT kill more mutants overall?"
\subsection{Comparison with prior work}
These results are consistent with the earlier findings of Jimenez et al. \cite{jimenez2018naturalness}: selection based on language models does not overwhelmingly outperform random selection - they found $A_{12} \approx 0.57-0.58$ (small) for n-gram naturalness, while we found a pooled $r\_rb=+0.350$ for the LLM. However, there is a noteworthy qualitative difference: our effect is stratified by code type ($r\_rb$ ranging from +0.028 in Codec to +0.571 in CSV and +0.436 in JacksonCore), whereas Jimenez's results were relatively uniformly low. This suggests that a modern LLM grasps semantics well enough to make a difference in complex code - an advantage that n-gram surface statistics lack. Compared with PMT approaches \cite{kim2022seshat, jain2023mutationbert, zhao2024soda}, our approach is completely zero-shot (training-free, requiring no historical kill data - whereas AL-PMT still needs 10\% \cite{borsi2025alpmt}); in return, there is no calibrated probability, only a raw score from 1-5 - which is the source of the ties analyzed above. Finally, the +4.6-point MS difference from mutant selection alone is far more modest than the improvement from mutant generation using LLMs (+32.99 points in real-bug detection \cite{wang2025comprehensive}) - which intuitively makes sense: generation opens up new space, whereas selection only optimizes within the available pool.
\subsection{Implications}
For practitioners: in a budget-constrained context, GPT-4o-mini is a safe and inexpensive choice - no worse than random selection in any project, clearly better on complex parser classes, and requiring no training or GPU infrastructure. However, for a codebase that already has a strong test suite (a high-MS architecture like Codec), investing in LLM-based selection is almost unprofitable - random selection is sufficient. For researchers: (1) the unit of analysis determines the conclusion - selection studies should report both group-weighted and mutant-weighted results and describe the structure of ties, in the same spirit as the pitfalls systematized by Delgado-Pérez et al.; (2) the raw 1-5 scale is a bottleneck - higher-resolution output (probabilities, pairwise ranking) could unlock performance currently being swallowed by ties.